\section{Discussion}
\label{sec:discussion}

\paragraph*{The envelope is generic; the rules add the classes.}
The main lesson of the production runs is a negative one. A long-lived demography alone, without any of the new
rules, drives 199 of 200 colonies to extinction: once natality stops, the colony simply ages out. Neither extinction
nor the rest of the envelope is a result of the rules. With crowding damage and the calibrated attraction, the
candidate without its four rules has a low peak with free space (the primary O4 passes in 200 of 200 runs, with a
peak at 0.45 of the damage-threshold capacity and 74\% of the cells empty), an end of natality and no rebound;
scored at the end of the plasticity window, its late-born cohorts also fail to acquire social competence (O11 in
0.755 of the runs, and 0.81 in the long-lived control), through accumulated crowding damage. Of the eight essential
criteria, five (O4, O5, O6, O7 and O10) pass in at least two of the four null models and carry no evidence for the
rules. What the rules add is the coexistence of crowded and isolated deteriorated classes (O13). It is absent in the
long-lived control, without crowding damage and without the rules, and present with refuges alone (1.00 in
\textsf{C1}$'$) or crowd aversion alone (0.92); full inheritance without social learning also produces it. Other
patterns hold by construction and only verify the implementation: the isolated class is persistent (O13p) because
the refuge capacity equals the isolation threshold, and transplanted late-phase animals never found a colony (O12)
because under R1 an adult cannot recover and conception scales with the social state of both partners. The
non-trivial part is upstream: the colony drives essentially every adult to $s\approx0$, and crowding damage alone
largely achieves it.

\paragraph*{A boom of one or two generations.}
The founders deteriorate by crowding during their own adult life; their offspring, about 70\% of all births, are
born below the deterioration threshold and learn from adults who have already deteriorated. The net reproductive
number falls from 3.8 in the founders to 0.40 in their daughters, so natality stops within about one generation.
The null models have a \emph{deeper} boom (mean generation of the births 1.59 without the four rules, 1.34 in
\textsf{C1}): the rules, R2 in particular, shorten the generational depth rather than create a long cascade. The
phases appear in Calhoun's order, but B and C are compressed and an artificial plateau precedes D. The ratio of the
deterioration time to the generation time, $\Gamma\approx1.1$ against about 3.5 in Universe~25 with a different
operational definition, is a heuristic comparison of time scales, not the cause of the shallow boom. Within this
transient, crowding damage is the trigger, social learning the amplifier, the plasticity window the ratchet, and the
refuges make the isolated class persistent; without R1 the refuges become pockets of low density where deteriorated
animals recover and breed again (the primary outcome falls to 0.21).

\paragraph*{The central limitation: demography and initial condition.}
The model has no mortality before senescence, and its initial condition is a synchronous cohort of 400 healthy
four-week-old animals. With a background hazard of 0.003 per week at all ages the primary outcome falls from 0.995
to 0.55 (0.06 at 0.008 per week), and with initial ages spread over 4--52 or 4--80 weeks it is 0.93 and 0.63. In
every case the only criterion that fails is O5, by its timing clause. Natality still ends with the population near
its maximum, but the deaths bring the start of the plateau forward (from 34 to 29 weeks at 0.003 per week) while the
end of natality hardly moves (40 to 43 weeks), so O5 becomes marginal, as in 83 of the 90 failing runs at 0.003 per
week and 73 of the 74 with ages over 4--80 weeks. O13, O13p and O11 pass in every run of these arms. The dependence
is thus on the timing of the end of natality relative to the start of the plateau, a generic criterion that
discriminates against no null model and rests on a post hoc choice of that start. It is neither evidence against
what the rules add nor for it, but the primary outcome of 0.995 overstates how robustly the full conjunction of
patterns is reproduced.

\paragraph*{Circularity, scope and minimality.}
The same patterns were used to \emph{choose} the rules and to \emph{report} them, with no pattern held out, so their
fulfilment shows that they are \emph{possible} under local rules, not that they were predicted. Global robustness
is limited and relative to the sampled box: over two 15-factor Latin hypercubes the mean primary outcome is 0.040
[0.026, 0.059], and 2.0\% of the points are Calhoun-like. With at most one lattice step per week the animals are
almost sessile on the scale of the deterioration, unlike Calhoun's mice, and the coexistence of classes is lost with
more movement sub-steps. Minimality is conditional on the thresholds: the minimal set R1, R2 and R7 (\textsf{C1}$'$)
gives the primary outcome in 0.835 [0.78, 0.88] of 200 runs, borderline, and in 0.33 with a class threshold of 15\%
instead of 10\%; crowd aversion contributes but is not necessary in the preregistered sense.

\paragraph*{What is not reproduced.}
From four founding pairs a colony spreads as a slow front, and no point of an exploratory founder plane is
Calhoun-like (primary outcome at most 0.37). The persistent isolated class is made of refugees of the early boom,
most of whom entered the refuges as pups, not of a late-born cohort like Calhoun's beautiful ones. Total infant
mortality and the female bias of the remnant are not reproduced, because conception and pup survival decline
together and there is no sex-specific mortality. The model of Gwizdałła and Duś \cite{GwizdallaDus2026} reproduces
the founder phase and the timing of the curve, which ours does not, with a daily step, fitted life tables and
deviation factors triggered by the total population. Ours shows that local rules, without population-level
switches, produce coexisting classes of deteriorated animals, one of them persistent, on top of an envelope that
crowding damage and the demography already give. Both agree that replicate variability is large and that Calhoun's
trajectory is one admissible outcome among many. Making the required patterns and their thresholds explicit and
testable run by run against null models is, in our view, the main methodological contribution of this work;
background mortality, gestation, a founder protocol with realistic mobility and held-out patterns are the natural
next steps. We make no claim about human populations.
